% Preamble for "Critical mean out-degrees of lattice neighbour graphs"
\IfFileExists{lmodern.sty}{\usepackage[T1]{fontenc}\usepackage{lmodern}}{}
\usepackage[british]{babel}
\usepackage{amsmath,amsthm,amssymb,mathtools}
\usepackage[margin=1.05in]{geometry}
\usepackage{microtype}
\setlength{\emergencystretch}{1em}
\usepackage{xcolor}
\definecolor{paperlink}{RGB}{0,0,114}
\usepackage{tikz}
\usetikzlibrary{arrows.meta,decorations.pathmorphing}
\usepackage{booktabs}
\usepackage{hyperref}
\hypersetup{colorlinks=true,linktoc=all,
  linkcolor=paperlink,citecolor=paperlink,urlcolor=paperlink,pdfstartview=Fit,
  pdftitle={Critical mean out-degrees of lattice neighbour graphs},
  pdfauthor={Benedikt Jahnel, in conversation with Claude Opus 5.5}}
\usepackage[capitalise,nameinlink]{cleveref}
\crefname{figure}{Figure}{Figures}
\Crefname{figure}{Figure}{Figures}
\usepackage{enumitem}
\setlist[itemize]{leftmargin=2em,itemsep=0.25em,topsep=0.35em}
\setlist[enumerate]{leftmargin=2.25em,itemsep=0.25em,topsep=0.35em}
\setlength{\parindent}{1.25em}
\setlength{\parskip}{0pt}
\usepackage{fancyhdr}
\pagestyle{fancy}
\fancyhf{}
\fancyhead[L]{\small Critical mean out-degrees of lattice neighbour graphs}
\fancyhead[R]{\small\thepage}
\setlength{\headheight}{14pt}
\numberwithin{equation}{section}

\usepackage{dsfont}
\newcommand{\1}{\mathds 1}

\theoremstyle{plain}
\newtheorem{theorem}[equation]{Theorem}
\newtheorem{proposition}[equation]{Proposition}
\newtheorem{lemma}[equation]{Lemma}
\newtheorem{corollary}[equation]{Corollary}
\newtheorem{conjecture}[equation]{Conjecture}
\theoremstyle{definition}
\newtheorem{definition}[equation]{Definition}
\theoremstyle{remark}
\newtheorem{remark}[equation]{Remark}

\newcommand{\Z}{\mathbb Z}
\newcommand{\N}{\mathbb N}
\renewcommand{\P}{\mathbb P}
\newcommand{\E}{\mathbb E}
\newcommand{\cD}{\mathcal D}
\newcommand{\cC}{\mathcal C}
\newcommand{\cG}{\mathcal G}
\newcommand{\cV}{\mathcal V}
\newcommand{\pc}{p_{\mathrm c}}
\newcommand{\kc}{k_{\mathrm c}}
\newcommand{\bond}{\mathrm{bond}}
\newcommand{\iid}{\mathrm{iid}}
\newcommand{\lra}{\longleftrightarrow}
\newcommand{\cP}{\mathcal P}
\newcommand{\cS}{\mathcal S}
\newcommand{\Piv}{\mathrm{Piv}}
\AtBeginDocument{\renewcommand{\c}{\mathrm c}}
\AtBeginDocument{\renewcommand{\d}{\mathrm d}}
\definecolor{colCH}{RGB}{0,0,200}
\definecolor{colBJ}{RGB}{200,0,0}
\definecolor{colLX}{RGB}{210,110,0}
\definecolor{colNew}{RGB}{0,140,0}
\newcommand{\cCH}[1]{\textcolor{colCH}{#1}}
\newcommand{\cBJ}[1]{\textcolor{colBJ}{#1}}
\newcommand{\cLX}[1]{\textcolor{colLX}{#1}}
\newcommand{\cNew}[1]{\textcolor{colNew}{#1}}
\usepackage{pgfplots}
\pgfplotsset{compat=1.17}
\usepgfplotslibrary{fillbetween}
\definecolor{serB}{HTML}{2A78D6}
\definecolor{serD}{HTML}{EB6834}
\definecolor{serU}{HTML}{1BAF7A}
\newcommand{\bc}{\beta_{\mathrm c}}
\usepackage{diagbox}
